% !TEX program = lualatex
% Auto-generated from YAML (v3) — 12: データサイエンス・AI（5）：深層学習

\documentclass[aspectratio=43,professionalfonts,handout,10pt]{beamer}
\usepackage{beamer_template}

\newcommand{\LectureNum}{12}
\newcommand{\LectureTitle}{データサイエンス・AI（5）：深層学習}
\newcommand{\CourseName}{情報リテラシー}
\newcommand{\TextPages}{-}
\newcommand{\NextTopic}{データサイエンス・AI（6）：生成AI}
\newcommand{\NextPages}{-}

\begin{document}
% --- Slide 1: title ---

\begin{frame}{\CourseName\quad 第\LectureNum 回「\LectureTitle 」}
  {\small \TermName\quad \TeacherName\quad ／\quad 教科書 \TextPages}

  \vfill

  \begin{block}{今日の流れ}
    \begin{enumerate}
      \item 深層学習とは
      \item ニューラルネットワークの基本構造
      \item 活性化関数・損失関数・最適化
      \item バックプロパゲーション
      \item 深層学習の応用例
    \end{enumerate}
  \end{block}
  \vfill

  \begin{block}{今日の目標}
    \begin{itemize}
      \item 深層学習について説明できる。
    \end{itemize}
  \end{block}
\end{frame}

% --- Slide 2: twocol ---

\begin{frame}{機械学習と深層学習の違い}
  \begin{columns}[T]
    \begin{column}{0.48\textwidth}
      \begin{block}{従来の機械学習}
        \small
        \begin{itemize}
          \item 特徴量を人間が設計する
          \item 例：顔認識 → 目の位置，鼻の形を手動で定義
          \item 専門知識が必要
          \item 少ないデータでも動くことがある
        \end{itemize}
      \end{block}
    \end{column}
    \begin{column}{0.48\textwidth}
      \begin{exampleblock}{深層学習（ディープラーニング）}
        \small
        \begin{itemize}
          \item 特徴量を自動で学習する
          \item 例：顔認識 → 画像から自動で重要な特徴を発見
          \item 大量のデータで高精度を実現
          \item 人間の脳の神経細胞（ニューロン）を模倣した構造
        \end{itemize}
      \end{exampleblock}
    \end{column}
  \end{columns}
\end{frame}

% --- Slide 3: terms ---

\begin{frame}{ニューラルネットワークの基本}
  \begin{description}[labelwidth=6em]
    \item[\kw{ノード}]
      計算を行う単位（ニューロンに相当）
    \item[\kw{重み（weight）}]
      ノード間の結合の強さを表す数値。学習によって調整される
    \item[\kw{層（layer）}]
      ノードの集まり。入力層・隠れ層・出力層から構成される
    \item[\kw{今日のモデル構造}]
      入力784 → 隠れ層128 → 隠れ層64 → 出力10（総パラメータ数：109,386個）
  \end{description}
\end{frame}

% --- Slide 4: twocol ---

\begin{frame}{活性化関数}
  \begin{columns}[T]
    \begin{column}{0.48\textwidth}
      \begin{block}{ReLU（レルー）}
        \small
        \begin{itemize}
          \item マイナスの値を0にする
          \item 隠れ層でよく使われる
          \item 計算が速い
          \item f(x) = max(0, x)
        \end{itemize}
      \end{block}
    \end{column}
    \begin{column}{0.48\textwidth}
      \begin{exampleblock}{Softmax（ソフトマックス）}
        \small
        \begin{itemize}
          \item 出力を確率に変換する（合計が1になる）
          \item 出力層でよく使われる
          \item 各クラスの確率を出力
          \item 例：数字「7」である確率 95\%
        \end{itemize}
      \end{exampleblock}
    \end{column}
  \end{columns}
\end{frame}

% --- Slide 5: points ---

\begin{frame}{学習の仕組み：バックプロパゲーション}
  \begin{block}{ポイント}
    \begin{itemize}
      \item 損失関数：予測と正解のズレを数値化。学習の目標は損失を小さくすること
      \item 最適化：損失を小さくするように重みを少しずつ自動調整する
      \item バックプロパゲーション（誤差逆伝播）：間違いを後ろから前に伝えて修正
      \item 流れ：予測する（前向き）→ 間違いを見つける → 原因を探る（後ろから前に）→ 重みを修正 → 繰り返す
      \item 結果：繰り返すほど精度が上がる
    \end{itemize}
  \end{block}
\end{frame}

% --- Slide 6: points ---

\begin{frame}{深層学習の応用例}
  \begin{block}{ポイント}
    \begin{itemize}
      \item 画像認識：顔認識（スマホロック解除）・医療画像診断（レントゲンから病変検出）・自動運転（歩行者検出）
      \item 自然言語処理：ChatGPT・Claude（会話AI）・機械翻訳・文章生成
      \item 音声認識：Siri・音声入力・リアルタイム翻訳
      \item 画像生成：Stable Diffusion・テキストから画像を生成
    \end{itemize}
  \end{block}
\end{frame}

% --- Slide 7: points ---

\begin{frame}{Colab実習：MNIST手書き数字認識}
  \begin{block}{ポイント}
    \begin{itemize}
      \item MNIST：手書き数字（0-9）の画像データセット。訓練60,000枚，テスト10,000枚
      \item セル1：データ読み込みと可視化
      \item セル2：モデル構築（総パラメータ数109,386個）
      \item セル3：学習と学習曲線（5エポックで約98\%の精度）
      \item セル4：予測と評価・混同行列（ヒートマップで可視化）
      \item セル5：正解/不正解の例（どの数字が間違えやすいか）
    \end{itemize}
  \end{block}

  \vfill

  \begin{exampleblock}{補足}
    \begin{itemize}
      \item エポック（epoch）：訓練データ全体を1回学習すること
      \item 学習の様子：Epoch1→92\%→Epoch2→96\%→Epoch3→97\%→Epoch4→98\%→Epoch5→98\%
    \end{itemize}
  \end{exampleblock}
\end{frame}

% --- チェックリスト ---

\begin{frame}{まとめ・チェックリスト}
  \begin{block}{自己チェック（できたら $\checkmark$ をつけよう）}
    \begin{itemize}
      \item[$\square$] ニューラルネットワークの基本構造（入力・隠れ・出力層）を理解した
      \item[$\square$] 活性化関数・損失関数・バックプロパゲーションの仕組みを学んだ
      \item[$\square$] 画像認識・自然言語処理など深層学習の応用例を把握した
    \end{itemize}
  \end{block}

  \vfill

  {\small 質問があれば授業後またはTeamsで受け付けます。}
\end{frame}

\end{document}
